\section{Verification}
\subsection{Part I: Introduction}
\subsubsection*{Software Quality}
\begin{defBox}[title=How to assure software quality? The big picture]
    Define a \red{quality assurance (QA) plan} (\enquote{\blue{Plan zur Qualitätssicherung}}) (must be integrated into the software development process)
    \begin{itemize}
        \item Constantly \red{assess} design \red{quality} (quantitative, qualitative)
        \begin{infoBox}
            Bewerten Sie die Designqualität kontinuierlich: sowohl quantitative als qualitative Methoden sollten angewendet werden, um die Qualität der Entwürfe zu messen
        \end{infoBox}
        \item Be proficient in applying time-tested \red{design principles}
        \begin{infoBox}
            Beherrschen Sie bewährte Entwurfsprinzipien: Nutzen Sie zeitlose Prinzipien des Softwaredesigns, um solide und wartbare Lösungen zu entwickeln
        \end{infoBox}
        \item Use tools and design techniques that help to achieve quality
        \item Use \red{design pattern} to learn from and reuse proven solutions to recurring problems
        \item Systematically \red{verify} correctness \& performance of design
        \item \red{Validate} fulfillment of requirements
    \end{itemize}
\end{defBox}
\subsubsection*{Validation vs. Verification}
\begin{defBox}[title=Validation]
    \red{Validation} ensures that a software system meets the \red{customer's expectation}
    \begin{center}
        \red{Are we building the right system?}
    \end{center}
    \begin{itemize}
        \item Is the implemented feature set as intended and complete?
        \item Does it fit into organization's workflow?
    \end{itemize}
\end{defBox}
\begin{defBox}
    Validation aligns \red{problem space} with \red{solution space}//
    It hinges on correctly performed \red{requirements analysis}//
    \blue{Cannot be done without involving customer / user}
\end{defBox}
\subsubsection*{Validation vs. Verification}
\begin{defBox}[title=Verification]
    \red{Verification} ensures correctness against \red{system specification}
    \begin{center}
        \red{Are we building the system right?}
    \end{center}
    \begin{itemize}
        \item Do the implemented functions work as specified?
        \item How does the system react to faults?
    \end{itemize}
\end{defBox}
\begin{defBox}
    Verification is performed in \red{solution spaces} by \red{manufacturer}
\end{defBox}
\begin{infoBox}[title=Verification]
    \begin{itemize}
        \item \fatsf{Definition:} Überprüfung, ob das Produkt korrekt gemäß seiner Spezifikation entwickelt wurde
        \item \fatsf{Ziel:} Sicherstellen, dass die Software den definierten Anforderungen, Regeln und Standards entspricht
        \item \fatsf{Fokus:} \enquote{Bauen wir das Produkt richtig?}
        \item \fatsf{Methoden:}
        \begin{itemize}
            \item Reviews (z.B. Code- oder Design-Reviews)
            \item Inspektionen
            \item Statische Tests
            \item Statische Analysen
        \end{itemize}
        \item \fatsf{Beispiel:} Prüfen, ob der Code korrekt implementiert wurde, basierend auf der Spezifikation oder ob das Design mit den Anforderungen übereinstimmt
    \end{itemize}
\end{infoBox}
\begin{infoBox}[title=Validation]
    \begin{itemize}
        \item \fatsf{Definition:} Überprüfung, ob das entwickelte Produkt tatsächlich die Bedürfnisse und Erwartungen der Nutzer erfüllt
        \item \fatsf{Ziel:} Sicherstellen, dass die Software \enquote{das richtige Produkt} ist, also den Zweck erfüllt
        \item \fatsf{Fokus:} \enquote{Bauen wir das richtige Produkt?}
        \item \fatsf{Methoden:}
        \begin{itemize}
            \item Funktionale Tests
            \item Benutzertests
            \item Usability-Tests
            \item Prototyping
            \item \fatsf{Beispiel:} Testen, ob ein Benutzer mit der Software eine bestimmte Aufgabe effizient lösen kann, oder ob die Funktionalitäten für den beabsichtigten Einsatz geeignet sind.
        \end{itemize}
    \end{itemize}
\end{infoBox}
\subsubsection*{Overview of Verification Techniques}
\begin{defBox}[title=Static]
    Static techniques do \red{not} require to actually execute code
    \begin{itemize}
        \item \red{Code review:} Applicable at \red{any} stage of development process
        \item \red{Static checking:} \red{Automated} analyses: type checks, bug finders, ...
        \item \red{Formale verification:} Ensures program satisfies a formal \red{property}
    \end{itemize}
\end{defBox}
\begin{defBox}[title=Dynamic]
    Dynamic techniques \red{do} require code to be executed
    \begin{itemize}
        \item \red{Testing:} Assert correct program behavior for selected \red{inputs}
        \item \red{Runtime monitoring:} Instrument program with safety/security \red{assertions} whose violation is indicated during execution
    \end{itemize}
\end{defBox}
\subsubsection*{What You Will Learn about Verification}
\begin{defBox}[title=In this Course:]
    \begin{itemize}
        \item \red{Code reviews} (\enquote{\blue{Code-Review}}) briefly (have to experience it to appreciate this technique)
        \item \red{Static checking} \& \red{Formal verification} briefly (contents of course \enquote{Formale Methode im Softwareentwurf})
        \item \red{Testing} discussed \blue{in detail} (by far the most importent verification technique in practice)
        \item \red{Runtime monitoring} (\enquote{\blue{Laufzeitbeobachtung}}) in passing (mostly relevant for specific application areas, for example, security)
    \end{itemize}
\end{defBox}
\subsection{Part II: Code Reviews}
\subsubsection*{Software Inspections: Code Reviews or Code Inspections}
\begin{defBox}[title=Definition (Code Review)]
    A \red{code review} is a structured inspection \red{process} of a piece of software, performed in a \red{team}, with the goal of finding possible \red{defects} (errors, poor code, deviations from conventions)
\end{defBox}
\blue{The goal is not to discuss broader design issues or to validate requirements}
\subsubsection*{Building Blocks of Code Reviews}
\begin{defBox}
    Prerequisite: up-to-date \red{code} and complete \red{specification} available Code needs not to be complete or executable
\end{defBox}
\begin{defBox}[title=Team Members Take Roles]
    Author/owner, inspector/reader, scribe, chair/moderator (minimal setup)
\end{defBox}
\begin{figure}[h]
    \centering
    \includegraphics[width=1\textwidth]{Verification Images/Building Blocks of Code Reviews.jpeg}
\end{figure}
\subsubsection*{Code Review: Check Lists}
\begin{defBox}[title=Code Reviews Driven by \fatsf{Check Lists}]
    Try to detect specific fault classes\\
    \blue{Data fault} All variables initialized?\\
    \blue{I/O fault} All input variables used?
\end{defBox}
\begin{defBox}
    Check lists must be \red{adapted} to project environment: Programming language, coding standards, used static analysis, etc.
\end{defBox}
\begin{defBox}[title=Speed?]
    \begin{itemize}
        \item Ca. 500 statements / hour at overview
        \item Ca. 100 statements / hour at inspection
    \end{itemize}
    Limit length of meeting to ca. 2 hours
\end{defBox}
\begin{figure}[h]
    \centering
    \includegraphics[width=1\textwidth]{Verification Images/Code Review: Check Lists.jpeg}
\end{figure}
\subsubsection*{Code Reviews: Advantages and Obstacles}
\begin{defBox}[title=Advantages of Code Reviews]
    \begin{itemize}
        \item Strong \red{empirical evidence} that they \blue{work} and \blue{save cost}
        \item Help to \red{distribute knowledge} of code among team members
        \item Can find defects \red{before} they show in tests (or even rolled out code)
        \item Can \red{improve} code quality (including documentation)
        \item \red{No executable} system required (find defects \blue{early})
    \end{itemize}
\end{defBox}
\begin{defBox}[title=Obstacles against Code Reviews]
    \begin{itemize}
        \item Requires \red{mature} personolities: team members may feel criticized
        \item Under time pressure: feeling that \red{not productive}
        \item Only works well when \red{properly} conducted
    \end{itemize}
\end{defBox}
\subsection{Part III: Software Testing}
\subsubsection*{Testing}
\enquote{Program testing can be used to show the presence of bugs, but never to show their absence!}
\begin{defBox}[title=Why is that so?]
    Non-trivial programs have infinitely many, or at least infeasibly many, different execution paths - it is impossible to test them all
\end{defBox}
\begin{defBox}
    Nevertheless, testing is an exceedingly \red{useful} and wide-spread verification method, even if sometimes only due to lack of better alternatives
\end{defBox}
\subsubsection*{What is a Program Test?}
\begin{defBox}[title=Constituents of a Program Test]
    \begin{itemize}
        \item The \red{system} or \red{implementation under test} (SUT, IUT): Die Software oder Komponente, die getestet wird
        \item The test \red{input(s)}
        \item The runtime environment, preparation, run, clean-up: \red{test harness} (\enquote{\blue{Test-Harnisch}}, \enquote{\blue{Testrahmen}}, \enquote{\blue{Gerüst}})
        \begin{infoBox}[title=Die Laufzeitumgebung]
            \begin{itemize}
                \item \fatsf{Definition:} Die Umgebung, in der die Tests ausgeführt werden, einschließlich Vorbereitung, Ausführung und Bereinigung
                \item \fatsf{Komponenten der Laufzeitumgebung:}
                \begin{itemize}
                    \item \fatsf{Test-Harnisch/Test-Rahmen/Test-Gerüst:}
                    \begin{itemize}
                        \item Werkzeuge oder Frameworks, die das SUI während der Tests unterstützen (z.B. Initialisierung von Daten, Simulation von Abhängigkeiten)
                        \item Beispiele: JUnit, pytest, Selenium
                    \end{itemize}
                    \item \fatsf{Vorbereitung (Setup):} Konfiguration der Testumgebung (z.B. Start von Datenbanken)
                    \item \fatsf{Bereinigung (Teardown):} Rücksetzen der Umgebung nach dem Test (z.B. Löschen von Testdaten)
                \end{itemize}
            \end{itemize}
        \end{infoBox}
        \item A test \red{verdict} (Testergebnis) (pass/fail) given by a \red{test oracle} (\enquote{\blue{Urteil}}. \enquote{\blue{Testorakel}})
    \end{itemize}
\end{defBox}
\subsubsection*{Test Levels}
\begin{figure}[h]
    \centering
    \includegraphics[width=1\textwidth]{Verification Images/Test Levels.jpeg}
\end{figure}
\subsubsection*{Test Plan}
\begin{defBox}[title=What is a Test Plan?]
    \blue{Detailed description of the testing process}\\
    \red{Document} intended to be read by humans
    \begin{itemize}
        \item \red{Work plan:} phases, schedule
        \item Testing \red{procedures}
        \item Explanation of the $\Rightarrow$ \red{test design}
        \item \blue{Test documentation as important as code documentation}
    \end{itemize}
\end{defBox}
\subsubsection*{Test Design}
\begin{enumerate}
    \item Identify and analyse the \red{responsibilities} of the IUT
    \begin{itemize}
        \item Use \red{pre- and postconditions} in use cases
        \item Use minimal and success \red{guarantees} in fully dressed use cases
        \item Analyze distribution of \red{responsibilities} in design
    \end{itemize}
    \item \red{Add} test cases based on
    \begin{itemize}
        \item Use Case, Design, Code \red{analysis}
        \item Suspicions, minimal success \red{guarantees}
        \item General \red{heuristics} (for example, domain/expression boundaries)
    \end{itemize}
    \item Determine for each test case how \red{verdict} is reached: Provide expected results, programmed or human \red{test oracle}
\end{enumerate}
\begin{defBox}
    Testing must be based on a \red{fault model} \\
    Exhaustive testing impossible: Assumptions needed, \red{where} faults are found, \red{which} need to be found
\end{defBox}
\subsubsection*{Test Automation Prospects}
\begin{defBox}
    Testing is a \red{very expensive} activity: Ca. 15-25\% of total development cost (higher estimated exist)
\end{defBox}
\begin{defBox}[title=Which Aspects of Testing can be \fatsf{Automated}?]
    \begin{enumerate}
        \item \red{Running the tests} (on nighly basis, after each change, etc.) (state-of-art, supported by build tools)
        \item Generating \red{test inputs} (increasingly done)
        \begin{enumerate}
            \item Code-driven
            \item Model-driven
            \item Data-driven
        \end{enumerate}
        \item Generating \red{test verdict}
        \begin{enumerate}
            \item Test \red{oracle} is small program, added to test harness (frequent)
            \item Oracle generated automatically from \red{formal specification} (rare)
        \end{enumerate}
    \end{enumerate}
\end{defBox}
\begin{infoBox}[title=Ausfürhren der Tests]
    \begin{itemize}
        \item \fatsf{Definition:} Automatisierung der Testausführung in regelmäßgien Abständen oder bei bestimmten Auslösern
        \item \fatsf{Beispiele:}
        \begin{itemize}
            \item Ausführen von Tests im Rahmen von \fatsf{Nighly Builds}
            \item Automatisches Testen nach jeder Codeänderung (z.B. in \fatsf{Continuous Integration/Continuous Deployment}-Pipelines)
        \end{itemize}
        \item \fatsf{Atueller Stand der Technik:}
        \begin{itemize}
            \item Weitgehend automatisiert und unterstützt durch moderne Build-Tools wie Jenkins, GitLab CI/CD, CircleCI oder Azure Pipelines
        \end{itemize}
        \item \fatsf{Vorteile:}
        \begin{itemize}
            \item Frühzeitiges Erkennen von Fehlern
            \item Konsistente Durchführung von Regressionstests
        \end{itemize}
    \end{itemize}
\end{infoBox}
\begin{infoBox}[title=Generierung von Testeingaben]
    \begin{itemize}
        \item \fatsf{Definition:} Automatisierung der Erstellung von Eingabedate, die für das Testen des Systems verwendet werden
        \item \fatsf{Ansätze:}
        \begin{enumerate}
            \item \fatsf{Code-basiert:}
            \begin{itemize}
                \item Eingaben werden durch Codeanalyse generiert (z.B. \fatsf{Coverage guided Fuzzing})
                \item Beispiel-Tools: AFL (American Fuzzy Lop) LibFuzzer
            \end{itemize}
            \item \fatsf{Modellbasiert:}
            \begin{itemize}
                \item Eingaben werden aus einem Modell generiert, das das Verhalten oder die Einschränkungen des Systems beschreibt
                \item Beispiel-Tools: QuickCheck, TLA+
            \end{itemize}
            \item \fatsf{Datengetrieben:}
            \begin{itemize}
                \item Eingaben werden aus externen Datensätzen oder Vorlagen abgeleitet
                \item Einsatzbereiche: Testen von APIs oder datenintensiven Anwendungen
            \end{itemize}
        \end{enumerate}
    \end{itemize}
\end{infoBox}
\begin{infoBox}[title=Generierung des Testergebnisses (Test Verdict)]
    \begin{itemize}
        \item \fatsf{Defintion:} Automatisierung der Bewertung, ob ein Test erfolgreich war oder fehlgeschlagen ist
        \item \fatsf{Ansätze:}
        \begin{enumerate}
            \item \fatsf{Testorakel als Teil des Test-Harnischs:}
            \begin{itemize}
                \item Ein \fatsf{kleines Programm} wird als Testorakel in den Test-Harnisch eingebettet, um die Testergebnisse zu bewerten (häufiger Ansatz)
                \item Beispiel: Vergleich von Soll- und Ist-Werten bei einer Funktion
            \end{itemize}
            \item \fatsf{Automatische Generierung des Orakels aus formalen Spezifikationen:}
            \begin{itemize}
                \item Das Testorakel wird aus einer formales Spezifikation abgeleitet (selten, da oft komplex)
                \item Einsatzbereiche: Formale Verifikationssysteme (z.B. TLA+, Z-Spezifikation)
            \end{itemize}
        \end{enumerate}
    \end{itemize}
\end{infoBox}
\subsubsection*{Test Automation: Running the tests}
\begin{defBox}[title=Tasks of Test Automation System]
    \begin{enumerate}
        \item Set up test environment
        \begin{itemize}
            \item Start servers (for example, database), establish network connection, register services
        \end{itemize}
        \item Start IUT
        \item Bring IUT into required pretest state configuration
        \begin{itemize}
            \item Load required data tables, create required objects
        \end{itemize}
        \item Set test inputs
        \item Evaluate output + record verdict
        \item Clean up environment
        \begin{itemize}
            \item Delete files, stop services, reset data
        \end{itemize}
    \end{enumerate}
\end{defBox}
\begin{defBox}
    \red{Manual} tasks remain when generating test automation scripts: Creating test \red{input}, writing test \red{oracle}
\end{defBox}
\subsubsection*{Test Goal}
\begin{defBox}[title=What can be Achieved by Testing?]
    Establish \red{sufficient} trust that system is (minimally) operational by exercising the interfaces between its parts
\end{defBox}
\red{Comprehensive testing} \blue{encompasses the following activities:}
\begin{enumerate}
    \item \red{Execute} test suite and pass \red{verdict} on each test case
    \item Use \red{coverage tool} to instrument the IUT
    \item Re-run test suite and judge reported \red{coverage}
    \item If necessary, \red{add} test cases to increase coverage
    \item \red{Stop} testing once coverage goal (Abdeckungsziel) is met \& all tests pass
\end{enumerate}
\begin{defBox}
    Remember: \red{Exhaustive} testing is not possible!
\end{defBox}
\subsubsection*{Test Input}
\begin{defBox}[title=Definition(Test(Data) Point)]
    A \red{test point} is a specific \red{value} for
    \begin{itemize}
        \item a test case \red{input} or
        \item a state variable
    \end{itemize}
    The test point is selected from a domain.\\
    The \red{domain} is the set of values that input or state variables can assume
\end{defBox}
\begin{infoBox}
    \begin{itemize}
        \item \fatsf{Testpunkt:} Ein bestimmter Wert für die Eingabe oder den Zustand während eines Tests
        \item \fatsf{Domain:} Die Gesamtheit aller möglichen Werte, die die Eingebe- oder Zustandsvariablen eines Systems annehmen können
    \end{itemize}
\end{infoBox}
\begin{defBox}[title=Heuristics for Test Point Selection]
    \begin{itemize}
        \item \red{Equivalence} classes (to be refined later) (only one test point for expected equivalent outcome)
        \item \red{Boundary} values (min/max of ordered domain, pivot of comparison)
        \item \red{Special} values (\fatsf{null}, other values with specific semantics)
    \end{itemize}
\end{defBox}
\subsubsection*{Test Case, Test Suite, Test Run}
\begin{defBox}[title=Definition (Test Case\, Test Suite)]
    A \red{test case} consists of:
    \begin{itemize}
        \item Pretest state of the IUT, 
        \item test point/conditions: test input,
        \item expected result.
    \end{itemize}
    A collection of test cases is called \red{test suite}
\end{defBox}
\begin{defBox}[title=Definition (Test Run)]
    A \red{test run} is the execution (with results) of a test suite on the IUT.\\
    A test whose results match the expected results gets the \blue{verdict} \red{pass}, otherwise the test \red{failed}
\end{defBox}
\subsubsection*{Test Driver, Fault, Failure}
\begin{defBox}[title=Definition(Test Driver\, Test Harness)]
    \begin{itemize}
        \item A \red{test driver (\enquote{\blue{Treiber}})} is a class or utility program that applies test cases to an IUT
        \item The \red{test harness} (\enquote{\blue{Gerüst}}) is a system of test drivers and other tools to support test execution
    \end{itemize}
\end{defBox}
\begin{defBox}[title=Definition (Fault\, Failure)]
    \begin{itemize}
        \item A software \red{fault} (\enquote{\blue{Fehler}}) is missing or incorrect code
        \item A \red{failure} (\enquote{\blue{Versagen}}) is the manifested inability of a system / component to perform a required function within specified limits
    \end{itemize}
\end{defBox}
\subsubsection*{A Popular Unit Test Framework: JUnit}
\begin{codeBlock}{}
    public class AccountTest{
        Account account;

        @BeforeEach
        public void setup(){
            account = new Account(100);
        }

        @AfterEach
        public void tearDown(){ account = null; }

        @Test
        public void successfulWithdrawTest(){
            boolean success = account.withdraw(50);
            assertTrue(success);
        }

        @Test
        public void unSuccessfulWithdrawTest(){
            assertFalse(account.withdraw(150));
        }
    }
\end{codeBlock}
\begin{defBox}[title=Test Framework: JUnit]
    \begin{itemize}
        \item Methods annotated with \red{@BeforeEach} establish pretest state
        \item Methods annotated with \red{@AfterEach} executed after test (perform cleanup)
        \item Methods annotated with \red{@Test} implement test: are executed
        \item \red{assertTrue, assertFalse, ...} evaluate result and record verdict
    \end{itemize}
\end{defBox}
\subsubsection*{Testing Case Study}
\begin{defBox}[title=Input]
    Read three integer values from the \red{command line}\\
    The three values represent the relative lengths of the sides of a triangle
\end{defBox}
\begin{defBox}[title=Output]
    Tells whether the triangle is
    \begin{itemize}
        \item \red{Scalene:} no two sides are equal
        \item \red{Isosceles:} exactly two sides are equal
        \item \red{Equilateral:} all sides are equal
    \end{itemize}
\end{defBox}
\begin{defBox}
    Create a \red{test suite} for this program that you consider to be \red{comprehensive}
\end{defBox}
\subsubsection*{Test Plan for Triangle Example (Excerpt)}
\begin{figure}[h]
    \centering
    \includegraphics[width=1\textwidth]{Verification Images/Test Plan for Triangle Example.jpeg}
\end{figure}
\subsection{Part IV: Test Coverage}
\subsubsection*{Test Coverage: When to Stop Testing?}
\begin{defBox}[title=Structural Criteria (Code Structure)]
    \red{Statement Coverage (SC)} each \red{statement} executed at least once\\
    \red{Basic Block Coverage (BBC)} each \red{basic block} executed at least once\\
    \red{Branch Coverage (BC)} each \red{outgoing edge} from a node in the control flow graph executed at least once\\
    \red{Path Coverage (PC)} each \red{path} in the control flow graph executed at least once
\end{defBox}
\begin{defBox}[title=Logic-based Criteria (Logical Case Distinctions)]
    \red{Condition Coverage (CC)} each \red{condition} evaluated as true and as false in at least one test run\\
    \red{Decision Coverage (DC)} each \red{decision} (guard) evaluated as true and as false in at least one test run\\
    \red{Modified Condition Decision Coverage (MCDC)} combines aspects of CC, DC \& independence test\\
    \red{Multiple-condition Coverage (MCC)} all true-false \red{combinations} of conditions exercised in at least one test run
\end{defBox}
\clearpage
\begin{defBox}[title=Structural Coverage]
    Structural Coverage is based on the \red{control flow graph} (CFG) of a program
\end{defBox}
\begin{figure}[h]
    \centering
    \includegraphics[width=0.75\textwidth]{Verification Images/Structural Converage and CFG.jpeg}
\end{figure}
\begin{defBox}[title=Definition]
    An \red{execution path} of a program is a \blue{complete} path through the CFG. \\
    A path that can actually be taken is called \red{feasible}.
\end{defBox}
\subsubsection*{Structural Coverage: Statement Coverage}
\begin{defBox}[title=Definition (Statement Coverage (SC) Criterion)]
    A test suite TS achieves \red{statement coverage} for a program p if:\\
    for any statement s in p there is a test in TS under which s is executed.
\end{defBox}
\begin{figure}[h]
    \centering
    \includegraphics[width=1\textwidth]{Verification Images/Statement Coverage.jpeg}
\end{figure}
\subsubsection*{Structural Coverage: Basic Block Coverage}
\begin{defBox}[title=Coverage Criterion for Byte Code]
    \red{Basic block coverage} very similar to statement coverage - main difference usually applied to machine or byte code
\end{defBox}
\begin{defBox}[title=Definition (Basic Block Reminder)]
    A \red{basic block} is a \red{maximal} sequence of \red{non-branching} program statements (or instructions) that are always executed \red{together} or not at all.\\
    Execution of a basic block starts with the \red{first statement}, only the \red{final} statement may branch or return (jump).
\end{defBox}
\begin{figure}[h]
    \centering
    \includegraphics[width=1\textwidth]{Verification Images/Basic Block Coverage.jpeg}
\end{figure}
\subsubsection*{Structural Coverage: Branch Coverage}
\begin{defBox}[title=Definition (Branch Coverage (BC))]
    A test suite TS achieves \red{branch coverage} for a program p if:\\
    \red{Each edge} in the CFG (control flow branch) is executed by some test
\end{defBox}
\begin{figure}[h]
    \centering
    \includegraphics[width=1\textwidth]{Verification Images/Branch Coverage 1.jpeg}
\end{figure}
\begin{figure}[h]
    \centering
    \includegraphics[width=1\textwidth]{Verification Images/Branch Coverage 2.jpeg}
\end{figure}
\subsubsection*{Structural Coverage: Path Coverage}
\begin{defBox}[title=Definition (Path Coverage (PC))]
    \red{Path coverage} is satisfied by a test suite TS if for every \red{execution path} ep through the CFG there is at least one test in TS causing ep to be taken
\end{defBox}
\begin{codeBlock}{}
    if (i >= 0){
        m = 1;
    }
    if (j < 0){
        m = 2;
    }
\end{codeBlock}
\begin{itemize}
    \item How many execution paths? 4
    \item Test suite with PC?\\
    TC $\overset{\wedge}{=}$ (i, j): \{(-1, 1), (-1, -1), (1, 1), (1, -1)\}
\end{itemize}
\blue{PC implies BC}
\begin{defBox}
    PC unachievable in practice: unbounded \# of paths or not all paths feasible
\end{defBox}
\subsubsection*{Logic Coverage: Condition vs. Decision}
\begin{figure}[h]
    \centering
    \includegraphics[width=1\textwidth]{Verification Images/Logic Coverage.jpeg}
\end{figure}
\begin{defBox}[title=Definition(Condition)]
    A \red{condition} is a Boolean code expression that cannot be further divided into Boolean subexpressions:
    \begin{itemize}
        \item Boolean variable or literal
        \item Comparision expression between variables or literals
        \item \enquote{instanceof} expression, etc.
    \end{itemize}
\end{defBox}
\begin{defBox}[title=Definition(Decision)]
    A \red{decision} is the Boolean expression constituting the \blue{guard} of a conditional or loop statement.
\end{defBox}
\subsubsection*{Logic Coverage: Decision Coverage}
\blue{Let p be a program and D(p) the set of all} \red{decisions} \blue{in p}
\begin{defBox}[title=Definition (Decision Coverage)]
    A test suite TS achieves \red{decision coverage} if every d $\in$ D(p) is evaluated at least once to false and once to true by tests in TS.\\
    (Assume d is not always false or true)
\end{defBox}
\begin{defBox}[title=Example]
    if (a >= 0 \& b < 0){ i = 0; }
    \begin{itemize}
        \item Test suite with decision coverage? TC $\overset{\wedge}{=}$ (a, b) : \{(-1, -1), (1, -1)\}
    \end{itemize}
\end{defBox}
\subsubsection*{Logic Coverage: Condition Coverage}
\blue{Let p be a program and C(p) the set of all} \red{occurrences of} \blue{conditions in D(p)}
\begin{defBox}[title=Definition(Condition Coverage)]
    A test suite TS achieves \red{condition coverage} if every c $\in$ C(p) is evaluated t least once to false and once to true by tests in TS.
\end{defBox}
\begin{defBox}[title=Example]
    if (a >= 0 \& b < 0) \{ i = 0; \}
    \begin{itemize}
        \item Test suite with condition coverage? TC $\overset{\wedge}{=}$ (a, b) : \{(-1, 1), (1, 1), (-1, -1)\}
        \item Test suite with condition coverage \red{and} decision coverage?\\
        TC $\overset{\wedge}{=}$ (a, b) : \{(-1, 1), (1, -1)\}
    \end{itemize}
\end{defBox}
\subsubsection*{Logic Coverage: Modified Condition Decision Coverage}
\begin{defBox}[title=Definition (Modified Condition-Decision Coverage (MCDC))]
    For one \red{occurrence} of condition c \red{inside} decision d,\\
    MCDC is satisfied by a test suite TS if 
    \begin{enumerate}
        \item it evaluated d at least twice,
        \begin{itemize}
            \item once where c evaluates to false,
            \item once where c evaluates to true,
        \end{itemize}
        \item d evaluates differently in both,
        \item all other conditions in d evaluate identically in both \red{or} these are not evaluated at all in at least one case.
    \end{enumerate}
    MCDC is satisfied by TS if it satisfies MCDC for all occurrences c $\in$ C(p).
\end{defBox}
\begin{defBox}
    MCDC often required when developing safety-critical software
\end{defBox}
\blue{\enquote{Each condition must have an influence on the evaluation of the decision and is independent of other conditions}}
\begin{defBox}[title=Example (MCDC)]
    Decision: \red{a>=0 \& (b<0 | c)}
    \begin{itemize}
        \item Does TS = \{(1, 1, true), (-1, 1, false)\}, TC $\overset{\wedge}{=}$ (a, b, c) satisfy MCDC for condition \blue{a>=0}?\\
        \red{No} (does not evaluate condition \blue{c} to same value in both test cases)
        \item Does TS = \{(1, 1, false), (-1, 1, false)\} satisfy MCDC for condition \blue{a>=0}?\\
        \red{No} (does not evaluate whole decision differently)
        \item Does TS = \{(1, 1, true), (-1, 1, true)\} satisfy MCDC for condition a>=0?\\
        \blue{Yes}
    \end{itemize}
\end{defBox}
\clearpage
\subsection{Part V: Test Automation \& Tool Support}
\subsubsection*{Automated Test Case Generation (ATCG)}
\begin{infoBox}
    Das Schreiben von Testfällen mit guter Abdeckung ist sehr zeitaufwendig. ATCG bietet zwei grundlegende Ansätze:
    \begin{enumerate}
        \item \fatsf{White-Box-Ansatz (code-basiert)}\\
        Hierbei wird der Code der zu testenden Einheit (IUT - "Implementation Under Test") analysiert, um eine möglichst hohe Abdeckung zu erreichen
        \begin{itemize}
            \item \fatsf{Syntaktische Ansätze:}
            \begin{itemize}
                \item Scannen nach Bedingungen und deren Auswertung
                \item Geeignet für logikbasierte Abdeckungsziele
                \item Kann Überapproximation enthalten (z.B. unerreichbarer Code oder nicht realisierbare Entscheidungen)
            \end{itemize}
            \item \fatsf{Symbolische Ausführung:}
            \begin{itemize}
                \item Entfaltung des Kontrollflussgraphen (CFG) mit symbolischen Werten
                \item Eignet sich auch für strukturelle Abdeckungskriterien
                \item Kann Unterapproximation beinhalten (z.B. nicht getesteter, aber erreichbarer Code)
            \end{itemize}
        \end{itemize}
        \item \fatsf{Black-Box-Ansatz}\\
        Hier wird entweder die Eingabedatenstruktur oder ein Modell des IUT analysiert, ohne direkten Block auf den Code
        \begin{itemize}
            \item \fatsf{Herausforderungen:}
            \begin{itemize}
                \item Bibliothekaufrufe, die schwer zu testen sind
                \item Irrelevante Testfälle, die keine zusätzlichen Erkenntnisse bringen
                \item Tiefe Schleifen oder Rekursionen, die schwer abzudecken sind
            \end{itemize}
        \end{itemize}
    \end{enumerate}
\end{infoBox}
\subsubsection*{Further Test Automation}
\begin{infoBox}[title=Test Coverage Recording (Testabdeckungs-Aufzeichnung)]
    Dieser Ansatz hilft nicht beim Schreiben von Testfällen, sondern dabei zu erkennen, wann genug getestet wurde.\\
    \fatsf{Schritte:}
    \begin{enumerate}
        \item \fatsf{Instrumentierung der IUT (Implementation Under Test)}
        \begin{itemize}
            \item Der Code wird mit zusätzlichen Anweisungen versehen, zm zu verfolgen welche Teile ausgeführt werden
        \end{itemize}
        \item \fatsf{Testausführung und Datenaufzeichnung}
        \begin{itemize}
            \item Während die Testsuite läuft, werden Abdeckungsdaten gesammelt
        \end{itemize}
        \item \fatsf{Analyse und Anzeige der Testergebnisse}
        \begin{itemize}
            \item Berichte, welche Codebereiche getestet wurden und welche nicht
        \end{itemize}
    \end{enumerate}
    \fatsf{Beispiel:}\\
    Ein Entwickler verwendet \fatsf{JaCoCo} für Java oder \fatsf{Coverage.py} für Python, um zu sehen, welche Codezeilen ausgeführt wurden
\end{infoBox}
\begin{infoBox}[title=Test Oracle Synthesis (Erstellung von Test-Orakeln)]
    Ein \fatsf{Test-Orakel} entscheidet, ob ein Test bestanden oder fehlgeschlagen ist\\
    \fatsf{Arten von Test-Orakeln:}
    \begin{itemize}
        \item \fatsf{Manuelles Orakel:}
        \begin{itemize}
            \item Ein Tester überprüft das Ergebnis $\rightarrow$ zeitaufwendig, fehleranfällig
        \end{itemize}
        \item \fatsf{Kodiertes Orakel (z.B. Assertions in JUnit):}
        \begin{itemize}
            \item Erfordert Expertenwissen, schwer zu pflegen
        \end{itemize}
    \end{itemize}
    \fatsf{Lösung:} Spezifikation in einer \fatsf{formalen Sprache} schreiben und daraus ein Test-Orakel generieren\\
    \fatsf{Herausforderungen:}
    \begin{itemize}
        \item Effiziente Generierung von Code aus einer \fatsf{deklarativen Spezifikation}
        \item Nicht alle Systeme lassen sich einfach in formale Regeln überführen
    \end{itemize}
\end{infoBox}
\subsubsection*{Automatic Static Verification Techniques}
\begin{infoBox}[title=Static Checking (Statische Überprüfung)]
    \fatsf{Wie funktioniert es?}
    \begin{itemize}
        \item Analysiert den \fatsf{Kontrollflussgraphen (CFG)} der IUT (Implementation under Test)
        \item Nutzt \fatsf{Constraint Solving} zur Erkennung von \fatsf{Laufzeitfehlern, Endlichkeit (Liveness) und Informationsfluss-Problemen}
        \item Vollautomatisch, kann aber \fatsf{zu viele false positives} erzeugen
        \item \fatsf{Polyzeit-Algorithmus}, d.h. skaliert gut für große Codebasen
    \end{itemize}
    Ein Tool wie \fatsf{SpotBugs} erkennt den Fehler \fatsf{ohne den Code auszuführen}
\end{infoBox}
\begin{infoBox}[title=Bug Finding (Fehlersuche durch Heuristiken \& Mustererkennung)]
    \fatsf{Wie funktioniert es?}
    \begin{itemize}
        \item Nutzt \fatsf{Pattern Matching \& Heuristiken} zur Erkennung typischer Bugs
        \item Sehr schnell, verarbeitet große Java-Codebasen effizient
        \item \fatsf{Über- und unter-approximiert} $\rightarrow$ kann Fehler übersehen oder zu viele Warnungen liefern
    \end{itemize}
    \fatsf{Java: SonarLint}
    \begin{itemize}
        \item \fatsf{Vortiel:} Findet typischer Fehler schnell
        \item \fatsf{Nachteil:} Kann wichtige Fehler übersehen (false negatives)
    \end{itemize}
\end{infoBox}
\subsubsection*{Automated Static Program Checking}
\begin{infoBox}[title=Kontrollflussanalyse (Control Flow Analysis\, CF)]
    \fatsf{Ziel:} Bestimmt, wie sich der Kontrollfluss im Program verhält
    \begin{figure}[h]
        \centering
        \includegraphics[width=1\textwidth]{Verification Images/Automated Static Program Checking.jpeg}
    \end{figure}
\end{infoBox}
\begin{infoBox}[title=Datenflussanalyse (Data Flow Analysis\, DF)]
    \fatsf{Ziel:} Analysiert, ob Werte von Variablen wirklich genutzt werden
    \begin{itemize}
        \item \fatsf{Lebensdaueranalyse (Liveness Analysis)}
        \begin{itemize}
            \item \fatsf{Ist x nach der Zuweisung genutzt?}
            \begin{itemize}
                \item Ja, x wird in z = x + 2; verwendet
            \end{itemize}
            \item \fatsf{Wird x = 5; je genutzt?}
            \begin{itemize}
                \item \fatsf{Nein}, weil x im if/else überschrieben wird \fatsf{$\rightarrow$ toter Code} 
            \end{itemize}
        \end{itemize}
        
    \end{itemize}
\end{infoBox}
\begin{infoBox}[title=Informationsflussanalyse (Information Flow Analysis\, IF)]
    \fatsf{Abhägigkeiten von z}
    \begin{itemize}
        \item z = x + 2;
        \item x hängt von if(y>0) ab $\rightarrow$ \fatsf{z ist abhängig von y}
    \end{itemize}
    \fatsf{Fazit:}
    \begin{itemize}
        \item Der Wert von z hängt direkt von y ab
        \item Falls y > 0, stammt x aus x = 1
        \item Falls y <= 0, stammt x aus x = 2
    \end{itemize}
\end{infoBox}
\begin{infoBox}[title=Pfadabhängigkeit (Path Analysis\, PA)]
    \fatsf{Ziel:} Bestimmt, welche Anfangswerte die Kontrollflüsse beeinflussen
    \begin{itemize}
        \item \fatsf{Welche Werte von y bestimmen den Pfad?}
        \begin{itemize}
            \item Falls y > 0 $\rightarrow$ linker Pfad (x = 1)
            \item Falls y <= 0 $\rightarrow$ rechter Pfad (x = 2)
        \end{itemize}
        \item \fatsf{Ergebnis:}\\
        y entscheidet, ob z = 3 (wenn y > 0) oder z = 4 (wenn y <= 0) ist
    \end{itemize}
\end{infoBox}
\subsubsection*{Bug Finding: SPOTBUGS}
\begin{defBox}[title=SPOTBUGS (spotbugs.github.io\, successor of FindBugs)]
    \begin{itemize}
        \item Bug detection in JAVA
        \item Finds \red{bug patterns} (idiom for error-prone code)
        \item Works at byte code level
    \end{itemize}
\end{defBox}
\begin{defBox}[title=Sources for Bug Patterns]
    \begin{itemize}
        \item Complex language features
        \item Misunderstood API methods or invariants
        \item Typos, wrong usage of operators
    \end{itemize}
\end{defBox}
\begin{defBox}
    False warnings: in practice < 50\% (good enough?)
\end{defBox}
\subsubsection*{SPOTBUGS: Some Bug Patterns}
\begin{infoBox}[title=Nutzloser Kontrollfluss zur nächsten Zeile (Useless Control Flow to next Line)]
    \begin{itemize}
        \item \fatsf{Problem:}
        \begin{itemize}
            \item Ein \fatsf{überflüssiges if-Statement} ohne Wirkung
            \item Der Code läuft \fatsf{unabhängig von der Bedingung} gleich weiter
        \end{itemize}
        \item \fatsf{Fehlerhafter Code:}
        \begin{codeBlock}{}
            if (args.length == 0);
                return;
        \end{codeBlock}
        \item \fatsf{SpotBugs-Warnung:}
        \begin{codeBlock}{}
            Useless control flow: control flow continues to same location regardless of branch taken
        \end{codeBlock}
    \end{itemize}
\end{infoBox}
\begin{infoBox}[title=Sinnloser Vergleich mit NaN (Doomed Test for Equality to NaN)]
    \begin{itemize}
        \item \fatsf{Problem:}
        \begin{itemize}
            \item Vergleich von double oder float mit NaN mittels == schlägt immer fehl
            \item NaN ist niemals gleich sich selbst (NaN == NaN) ist false
        \end{itemize}
        \item \fatsf{Fehlerhafter Code:}
        \begin{codeBlock}{}
            double x = Double.NaN;
            if (x == Double.NaN){ // Immer false!
                System.out.println("x ist NaN");
            }
        \end{codeBlock}
        \item \fatsf{SpotBugs-Warnung:}
        \begin{codeBlock}{}
            Doomed test for equality to NaN: 'x == Double.NaN' will always evaluate to false
        \end{codeBlock}
        \item \fatsf{Lösung:}
        \begin{codeBlock}{}
            if (Double.isNaN(x)){ // Korrekte Methode zur NaN-Prüfung
                System.out.println("x ist NaN");
            }
        \end{codeBlock}
    \end{itemize}
\end{infoBox}
\subsubsection*{Formal Verification}
\begin{infoBox}
    \fatsf{Formale Verfikation} ist eine Methode zur mathematischen \fatsf{Beweisführung}, um sicherzustellen, dass ein Programm seine \fatsf{formale Spezifikation} erfüllt
\end{infoBox}
\begin{infoBox}[title=Warum formale Verifikation?]
    \begin{itemize}
        \item \fatsf{Mathematische Strenge:} Basiert auf \fatsf{Logik und Mengenlehre}
        \item \fatsf{Starke Garantien:} Wenn eine Eigenschaft bewiesen wurde, \fatsf{gibt es keine Gegenbeispiele}
        \item \fatsf{Nicht immer vollständig:} Einige wahre Eigenschaften können \fatsf{nicht bewiesen werden}
    \end{itemize}
\end{infoBox}
\begin{infoBox}[title=Arten der formalen Verifikation]
    \begin{enumerate}
        \item \fatsf{Modellprüfung (Model Checking)}
        \begin{itemize}
            \item \fatsf{Prinzip:} Das Programm wird als \fatsf{endliher Zustandsraum} modelliert, und ein Algorithmus prüft \fatsf{alle möglichen Zustände}
            \item \fatsf{Beispiel-Tools:}
            \begin{itemize}
                \item \fatsf{CBMC} - Model Checking für C/C++
                \item \fatsf{Java PathFinder (JPF)} - Model Checking für Java
                \item \fatsf{Sping} - Verifikation von Parallelprogrammen
                \item Spin
                \item SDV/SLAM, ...
            \end{itemize}
            \item \fatsf{Nachteil: Zustandsexplosion} - Wenn das Programm viele Variablen hat, kann die Anzahl der möglichen Zustände extrem groß werden
        \end{itemize}
        \item \fatsf{Deduktive Verifikation (Deductive Verification)}
        \begin{itemize}
            \item \fatsf{Prinzip:} Eigenschaften des Programms werden \fatsf{als mathematische Beweise} formuliert
            \item \fatsf{Beispiel-Tools:}
            \begin{itemize}
                \item \fatsf{KeY} - Beweist Korrektheit von Java-Programmen
                \item \fatsf{KIV} - Beweist funktionale Korrektheit von Algorithmen
                \item \fatsf{VeriFast} - Verifikation von Multi-Threading-Programmen
                \item VeriFast, jStar, VerCors, Viper, ...
            \end{itemize}
            \item \fatsf{Vorteil:} Funktioniert auch für \fatsf{unendliche Zustandsräume}
            \item \fatsf{Nachteil:} Erfordert oft \fatsf{manuelle Annotationen} im Code, um den Beweis zu führen
        \end{itemize}
    \end{enumerate}
\end{infoBox}
\subsubsection*{Formal Verification: Architecture}
\begin{figure}[h]
    \centering
    \includegraphics[width=1\textwidth]{Verification Images/Formal Verification: Architecture.jpeg}
\end{figure}
\begin{infoBox}[title=Prozess der Verifikation]
    \begin{enumerate}
        \item \fatsf{Formale Spezifikation:}
        \begin{itemize}
            \item Definiert, \fatsf{was} ein Programm tun soll (z.B. \fatsf{Vorbedingungen, Nachbedingungen})
            \item \fatsf{Program Verification System:}
            \begin{itemize}
                \item Nutzt \fatsf{Beweisregeln,} um zu zeigen: \enquote{Die Implementierung erfüllt die Spezifikation}
            \end{itemize}
            \item \fatsf{Automatische Beweisführung:}
            \begin{itemize}
                \item Computerunterstützt (z.B. mit \fatsf{KeY} für Java)
            \end{itemize}
        \end{itemize}
    \end{enumerate}
\end{infoBox}
\subsubsection*{Model Checking: SDV/SLAM - Static Driver Verification}
\begin{infoBox}
    \fatsf{SDV (Static Driver Verifier) / SLAM} ist ein \fatsf{formales Verifikations-Tool} zur \fatsf{Fehlersuche in Windows-Gerätetreibern} (C-Programme)
\end{infoBox}
\begin{infoBox}[title=Funktionsweise von SDV/SLAM]
    \begin{itemize}
        \item Findet alle Fehler mit Bezug auf API-Nutzungsregeln
        \item Analysiert alle möglichen Codepfade durch Model Checking
        \item Definiert temporale Sicherheitsregeln für API-Nutzung
        \item Verwendet CEGAR (Conterexample-Guided Abstraction Refinement)
        \item Nutzt symbolische Ausführung \& Model Checking 
        \item Sehr geringe False-Positive-Rate (4\%)
    \end{itemize}
\end{infoBox}
\subsubsection*{Specifications as Contracts}
\begin{infoBox}[title=Spezifikationen als Verträge (\enquote{Design by Contract})]
    \begin{itemize}
        \item \fatsf{Grundidee:} Methoden in der objektorientierten Programmierung werden durch \fatsf{Verträge} beschrieben
        \item \fatsf{Vertrag = Spezifikation, die Bedingungen für Korrektheit definiert}
        \item \fatsf{Zwei Parteien:}
        \begin{itemize}
            \item \fatsf{Caller (Aufrufer)} $\rightarrow$ Muss \fatsf{Vorbedingungen} erfüllen
            \item \fatsf{Callee (Methode/Implementierung)} $\rightarrow$ Garantiert \fatsf{Nachbedingungen,} falls Vorbedingungen erfüllt sind
        \end{itemize}
    \end{itemize}
\end{infoBox}
\begin{infoBox}[title=\enquote{Design by Contract} (Meyer\, 1992\, Eiffel)]
    \begin{itemize}
        \item Methode garantiert ein bestimmtes Verhalten $\rightarrow$ Nur wenn der Caller die Voraussetzungen erfüllt!
        \item Verhindert unklare Verantwortlichkeiten zwischen Entwicklern
    \end{itemize}
    \fatsf{Drei Bestandteile eines Vertrags:}
    \begin{itemize}
        \item \fatsf{Precondition (Vorbedingung):} Bedingungen, die erfüllt sein müssen, bevor die Methode aufgerufen wird
        \item \fatsf{Postcondition (Nachbedingung):} Bedingungen, die nach der Methode garantiert erfüllt sind
        \item \fatsf{Invariant:} Bedingungen, die für eine Klasse jederzeit gültig sind
    \end{itemize}
\end{infoBox}
\begin{infoBox}[title=Beispiel mit JML (Java Modeling Language)]
    \begin{itemize}
        \item Vertrag für eine withdraw-Methode in einer BankAccount-Klasse:
        \begin{codeBlock}{}
            public class BankAccount {
                private int balance;
                
                /*@ requires amount > 0 && amount <= balance;
                  @ ensures balance == \old(balance) - amount;
                  @*/
                public void withdraw(int amount) {
                    balance -= amount;
                }
            }
        \end{codeBlock}
        \item \fatsf{Erklärung:}
        \begin{itemize}
            \item requires amount > 0 \&\& amount <= balance;\\
            $\rightarrow$ \fatsf{Vorbedingung:} Betrag muss positiv und nicht größer als das Guthaben sein
            \item ensures balance == \textbackslash old(balance) - amount;\\
            $\rightarrow$ \fatsf{Nachbedingung:} Nach dem Aufruf ist das Guthaben um amount reduziert
        \end{itemize}
    \end{itemize}
\end{infoBox}
\subsubsection*{Java Modelling Language (JML)}
JML is a \blue{contract-based} \red{specification language} tailored to \red{Java}, a \red{behavioral interface sspecification language} (BISL)
\begin{defBox}[title=General JML Philosophy]
    Integrate
    \begin{itemize}
        \item JML specification and
        \item Java implementation
    \end{itemize}
    within \red{a single language} (\enquote{single-tier approach})
\end{defBox}
$\Rightarrow$ JML is not external to Java, but an \red{extension} of Java
\begin{defBox}
    JML \red{is} Java + \red{First-Order Logic} + \red{Contracts} + \red{Invariants} + more...
\end{defBox}
\subsubsection*{Specifying Binary Search in JML}
\begin{codeBlock}{}
    private int[]a;
    /*@ private normal_behavior
        requires 0 <= low <= up <= a.length;
        requires (\forall int x, y; 0 <= x < y < a.length; a[x] <= a[y]);
        ensures \result == -1 || low <= \result < up;
        ensures (\exists int idx; low <= idx < up; a[idx] == v) ? \result >= low && a[\result] == v : \result == -1;
        assignable \nothing;
        measured_by up - low;
    @*/
    private int binSearch (int v, int low, int up);
\end{codeBlock}
\begin{defBox}[title=Natural Language Specification]
    If the caller guarantees that
    \begin{enumerate}[label=\roman*)]
        \item low is less-or-equal than up, both are in the bounds of a (including a.length)
        \item a is not null (Default in JML Implicit precondition \fatsf{requires a != null})
        \item a is sorted
    \end{enumerate}
    then the method guarantees that
    \begin{enumerate}[label=\roman*)]
        \item the result is -1 or in [low, up)
        \item if v is in a then an index of v in a is returned, else -1 is returned
        \item it terminates w/o an exception
        \item the heap is not modified
        \item up - low is a termination witness
    \end{enumerate}
\end{defBox}
\subsubsection*{Working Principle of Deductive Verification}
\begin{infoBox}
    Deduktive Verifikation nutzt \fatsf{mathematische Beweise} zur Sicherstellung der Korrektheit eines Programms anhand seiner \fatsf{formalen Spezifikation}
\end{infoBox}
\begin{infoBox}[title=Path Exploration (Wege-Exploration)]
    \fatsf{Symbolische Ausführung} wird verwendet, um alle möglichen Pfade im \fatsf{Kontrollflussgraphen (CFG)} eines \fatsf{linearen Programms} (ohne Sprünge, Schleifen oder Methodenaufrufe) zu untersuchen
    \begin{itemize}
        \item \fatsf{Symbolische Werte:} Statt konkrete Eingaben zu testen, werden \fatsf{Variablen als Symbole} behandelt
        \item \fatsf{Endliche Anzahl von Pfaden:} Bei linearem Code gibt es \fatsf{nur eine begrenzte Anzahl von Ausfürhungspfaden}
    \end{itemize}
    Beispiel:
    \begin{codeBlock}{}
        int foo(int x) {
            if (x > 0) return x + 1;
            else return x - 1;
        }
    \end{codeBlock}
    \fatsf{Mögliche Pfade:}
    \begin{enumerate}
        \item \fatsf{x > 0} $\rightarrow$ return x + 1
        \item \fatsf{x <= 0} $\rightarrow$ return x - 1
    \end{enumerate}
    Durch symbolische Werte (x = $\alpha$, x = $\beta$) kann die Methode \fatsf{für alle Eingaben} formal verifiziert werden
\end{infoBox}
\begin{infoBox}[title=Umgang mit Schleifen - Loop Invariants]
    \begin{itemize}
        \item \fatsf{Problem:} Schleifen erzeugen unendlich viele mögliche Pfade
        \item \fatsf{Lösung: Schleifeninvarianten} approximieren die Wirkung einer Schleife statt sie unendlich auszuführen
    \end{itemize}
    \fatsf{Beispiel für eine Schleifeninvariante}
    \begin{codeBlock}{}
        int sum(int n) {
        int s = 0;
        int i = 1;
        while (i <= n) {
            s = s + i;
            i++;
        }
        return s;
    }
    \end{codeBlock}
\end{infoBox}
\fatsf{Schleifeninvariante:}
\enquote{Am Anfang jeder Iteration gilt: s ist die Summe der Zahlen 1 bis i-1}
\begin{itemize}
    \item Diese Regel bleibt für \fatsf{alle} iterationen gültig
    \item Das Verifikationssystem muss nur \fatsf{die Invariante beweisen}, nicht jede Iteration einzeln
\end{itemize}
\begin{infoBox}[title=Schlüssel zur Skalierbarkeit: Methodenverträge statt vollständiger Ausführung]
    \begin{itemize}
        \item Anstatt Methodenaufrufe vollständig zu symbolisch auszuführen, werden sie \fatsf{durch ihre Verträge ersetzt}
        \item Dies reduziert die Anzahl der analysierten Pfade drastisch
    \end{itemize}
    \fatsf{Beispiel:}
    \begin{codeBlock}{}
        //@ requires x > 0;
        //@ ensures \result == x + 1;
        int increment(int x) {
            return x + 1;
        }
        
        void test(int y) {
            int z = increment(y);
        }
    \end{codeBlock}
    \fatsf{Statt increment(y) symbolisch auszuführen}, wird sein Vertrag genutzt:
    \begin{itemize}
        \item y > 0 garantiert als Vorbedingung
        \item z == y + 1 ist direkt ableitbar
    \end{itemize}
    \fatsf{Ergebnis:} Statt alle Pfade innerhalb von increment() zu analysieren, wird der \fatsf{Vertrag direkt als Beweisregel verwendet!}
\end{infoBox}
\subsubsection*{Symbolic Program Execution}
\begin{figure}[h]
    \centering
    \includegraphics[width=1\textwidth]{Verification Images/Symbolic Program Execution.jpeg}
\end{figure}
\clearpage
\subsubsection*{Dynamic Verification: Runtime Monitoring of Contract Adherence}
\begin{table}[h!]
    \centering
    \renewcommand{\arraystretch}{1.5} % Adjust row height
    \begin{tabular}{|c|c|c|}
        \hline
        \fatsf{Merkmal} & \fatsf{Statische Verifikation (Deduktiv)} & \fatsf{Dynamische Verifikation (Runtime Monitoring)} \\
        \hline
        \fatsf{Wann erfolgt die Prüfung?} & Vor der Ausfürhung (Compile-Zeit) & Während der Programmausführung \\
        \hline
        \fatsf{Wie?} & Mathematischer Beweis & Laufzeit-Überprüfung mit assert \\
        \hline
        \fatsf{Erkennt Fehler?} & Für alle möglichen Eingaben & Nur für ausgeführte Fälle \\
        \hline
        \fatsf{Skalierbarkeit?} & Anspruchsvoll, aber vollständig & Einfacher, aber begrenzte Abdeckung \\
        \hline
    \end{tabular}
\end{table}
\begin{infoBox}[title=Grenzen der dynamischen Verifikation]
    \begin{itemize}
        \item \fatsf{Keine vollständige Abdeckung}
        \begin{itemize}
            \item Nur die \fatsf{ausgeführten Fälle} werden geprüft. Falls ein Codepfad nie durchlaufen wird, wird er \fatsf{nicht getestet}
        \end{itemize}
        \item \fatsf{Kein Beweis für alle Eingaben}
        \begin{itemize}
            \item Ein erfolgreicher Lauf bedeutet nicht, dass der Code für \fatsf{alle möglichen Eingaben} korrekt ist
        \end{itemize}
        \item \fatsf{Gut für Debugging \& Testfälle}
        \begin{itemize}
            \item Kann helfen, Fehler frühzeitig zu finden, insbesondere in der Entwicklung \& Testphase
        \end{itemize}
    \end{itemize}
\end{infoBox}
\begin{infoBox}[title=Kombination von Statischer und Dynamischer Verifikation]
    Die Kombination von \fatsf{statischer (deduktiver) Verifikation} und \fatsf{dynamischer Verifikation (Laufzeitüberwachung)} kann helfen, \fatsf{maximale Sicherheit} bei der Programmkorrektheit zu erreichen
    \begin{itemize}
        \item \fatsf{Statische Verifikation} (z.B. mit KeY oder JML) beweist mathematisch, dass ein Codeabschnitt für \fatsf{alle möglichen Eingaben} korrekt ist
        \item \fatsf{Dynamische Verifikation} (mit assert) stellt sicher, dass auch zur \fatsf{Laufzeit} keine unvorhergesehenen Fehler auftreten
    \end{itemize}
\end{infoBox}
\subsubsection*{Verification: Summary}
\begin{itemize}
    \item \blue{Testing} is by far the most important and common \red{verification approach}
    \item \red{Automated test case execution} is standard (and essential)
    \item \red{Automated test generation} tools require expertise, lack robustness, but:
    \begin{itemize}
        \item Potentially huge cost savings
        \item Test suites with coverage guarantees
        \item Several industrial products availale, growing interest
    \end{itemize}
    \item \red{Code reviews} often under-appreciated, but effective
    \item Fully automatic verification tools include \red{bug finders}, \red{static checkers}
    \begin{itemize}
        \item Main problem: high number of \red{false positives}
    \end{itemize}
    \item \red{Formal verification} is mainly done
    \begin{itemize}
        \item in \red{specific} application scenarios (domain-specific, hard-coded rules)
        \item in \red{safety-critical} systems (effort to \red{supply formal specification} is justified)
        \item for high-impact code, such as \red{libraries} (dto.)
    \end{itemize}
\end{itemize}